% Generated draft: requires mathematical and source review.
\documentclass[11pt]{article}
\usepackage[margin=25mm]{geometry}
\usepackage{fontspec}
\setmainfont{Latin Modern Roman}
\usepackage{amsmath,amssymb,mathtools,unicode-math}
\setmathfont{Latin Modern Math}
\usepackage{enumitem,needspace,xurl,hyperref}
\setlength{\parindent}{0pt}
\setlength{\parskip}{5pt}
\setlength{\emergencystretch}{3em}
\newcommand{\N}{\mathbb{N}}
\newcommand{\Z}{\mathbb{Z}}
\newcommand{\Q}{\mathbb{Q}}
\newcommand{\R}{\mathbb{R}}
\newcommand{\C}{\mathbb{C}}
\newcommand{\F}{\mathbb{F}}
\newcommand{\A}{\mathbb{A}}
\newcommand{\PP}{\mathbb{P}}
\newcommand{\E}{\mathbb{E}}
\newcommand{\eps}{\varepsilon}
\newcommand{\dd}{\,\mathrm{d}}
\begin{document}
\section*{UnsolvedMath ID: 30006293}
\textbf{Source problem code:} catalogue entry
\subsection{Definitions and mathematical statement}
A selective-decay variational principle describes the selection of flow states through decay of a designated functional subject to specified retained constraints. A forced-dissipative flow includes forcing and dissipation. Topography denotes the spatially varying underlying relief, and a topographically pinned vortex denotes a vortex whose position is tied to that relief. These descriptions identify the intended concepts; the supplied extract does not specify the governing equations, functionals, constraints, or a mathematical criterion for pinning.

The supplied problem asks for the following extensions and investigations:
\begin{enumerate}
\item Extend the selective-decay variational principles referred to in the source to out-of-equilibrium, forced-dissipative flows. Specify the applicable equations, forcing and dissipation, admissible states, variational functional, constraints, and sense in which the resulting principle predicts the flow.
\item Investigate whether the observed agreement between the referenced theory and numerical simulations persists for more complex topography. The source explicitly states that the reported simulations used single-scale topography, whereas the referenced theory holds for any topography. This assertion is retained as a source qualification; the extract does not supply the theory or its hypotheses, so it does not determine the mathematical meaning of ``any topography.''
\item Explain the physical origin of vortices pinned to the topography and extend the theory to predict their strengths and locations. The notions of vortex strength, location, and successful prediction must be specified for the chosen flow model.
\end{enumerate}
A fully quantified mathematical statement cannot be recovered from this extract: the spatial domain, boundary conditions, flow model, regularity assumptions, admissible topographies, forcing and dissipation regimes, and prediction criteria are unspecified. The questions above therefore remain a statement of the source's research targets, requiring those qualifications before a precise theorem or conjecture can be formulated.
\subsection{Short English statement}
Develop selective-decay variational principles for forced-dissipative flows, investigate whether the reported agreement with simulations persists for complex topography, and explain and predict the strengths and locations of topographically pinned vortices. The source reports simulations with single-scale topography and claims broader applicability of the referenced theory, but supplies insufficient model details for a fully quantified formulation. Its partially solved label and weak-forcing/weak-damping assessment are supplied dated curation, not independently verified conclusions.
\subsection{Sources}
\begin{itemize}
\item Local UnsolvedMath registry: supplied complete record, permanent ID \(30006293\), title \emph{Selective Decay in Forced-Dissipative Flows}; exact source problem code \texttt{catalogue entry}. The record assigns the proposed year 2025 and supplies no individual proposer.
\item \emph{Mathematical Advances in Geophysical Fluid Dynamics} (2025), Oberwolfach Reports, pp.\ 1215--1270, DOI 10.4171/owr/2025/24. The supplied passage includes the page marker 1261. Original report: \url{https://doi.org/10.4171/owr/2025/24}. The registry describes Oberwolfach Reports as documenting mathematical workshops and their open-problem discussions.
\item Supplied literature reference attributed to Gallet, ``two-dimensional turbulence above topography'': \url{https://www.cambridge.org/core/journals/journal-of-fluid-mechanics/article/twodimensional-turbulence-above-topography-condensation-transition-and-selection-of-minimum-enstrophy-solutions/654AD74779542F674C5A51FE5F82F28}. The supplied record does not provide complete bibliographic metadata.
\item Literature assessment dated exactly \texttt{2026-08-22}: the supplied curation labels the problem partially solved and reports a perturbative selective-decay extension for weak forcing and damping, with numerical validation of bifurcation predictions, while stating that a general forced-dissipative theory is not established. This is supplied dated curation, not an independently verified literature assessment. Its listed sources are the Gallet reference and the original Oberwolfach report linked above; the record explicitly calls for independent verification of time-sensitive status and cited claims.
\end{itemize}
\end{document}
